\begin{table}[htbp]
\centering
\footnotesize
\caption{Problema D, meia-lua (openai/gpt-4o-mini): protocolo da Fase A em R$^2$, R$^3$ e R$^4$.}
\label{tab:T3_D_gpt4mini}
\begin{tabular}{lrrrrlrrlr}
\toprule
 &  &  & \multicolumn{3}{c}{Perceptron} & \multicolumn{3}{c}{NNLS} &  \\ \cmidrule(lr){4-6} \cmidrule(lr){7-9}
Regime & Acur.\ LLM vs real & Partição LLM (n0/n1) & Fid.\ (vs LLM) & Acur.\ vs real & $\hat{W}$ unit. & Fid.\ (vs LLM) & Acur.\ vs real & $\hat{W}$ unit. & $\Delta$ fid.\ vs R$^2$ \\
\midrule
R$^2$ — zero-shot† & 50,7 & 149/1 & 66,7 & 63,3 (62,7–64,0) & [0,34; 0,94] & 68,0 (66,7–69,3) & 60,7 (60,0–61,3) & [0,37; 0,93] & +0,0 \\
R$^2$ — 4-shot† & 81,4 (68,5–91,8) & 83/63 & 89,1 (75,3–95,9) & 81,0 (76,0–85,3) & [0,48; 0,82] & 88,7 (77,4–96,6) & 83,3 (78,7–86,7) & [0,22; 0,97] & +0,0 \\
R$^2$ — 10-shot† & 89,0 (81,4–97,1) & 72/68 & 91,4 (80,0–97,1) & 83,9 (77,3–86,0) & [0,31; 0,94] & 91,9 (80,0–97,1) & 84,7 (81,3–86,0) & [0,27; 0,96] & +0,0 \\
R$^2$ — 20-shot† & 90,4 (86,2–95,4) & 68/62 & 92,6 (91,5–94,6) & 84,9 (82,7–87,3) & [0,29; 0,95] & 92,8 (90,0–96,2) & 85,7 (84,7–86,7) & [0,26; 0,97] & +0,0 \\
R$^2$ — 40-shot† & 91,9 (87,3–96,4) & 57/53 & 91,3 (88,2–92,7) & 86,1 (84,7–88,0) & [0,25; 0,97] & 91,5 (88,2–93,6) & 85,9 (84,7–86,7) & [0,23; 0,97] & +0,0 \\
R$^2$ — W ajustado ao rótulo real (teto)† & — & — & — & 85,6 (84,7–86,7) & [0,18; 0,98] & — & 86,7 & [0,22; 0,97] & — \\
\midrule
R$^3$ — zero-shot† & 50,7 & 149/1 & 71,0 (70,7–71,3) & 63,0 (62,7–63,3) & [0,21; 0,61; 0,77] & 81,0 (79,3–82,7) & 58,3 (57,3–59,3) & [0,06; 0,15; 0,99] & +4,3 \\
R$^3$ — 4-shot† & 81,4 (68,5–91,8) & 83/63 & 89,8 (78,8–95,9) & 81,9 (80,0–84,7) & [0,33; 0,63; 0,67] & 90,9 (77,4–99,3) & 82,5 (78,7–85,3) & [0,13; 0,60; 0,53] & +0,7 \\
R$^3$ — 10-shot† & 89,0 (81,4–97,1) & 72/68 & 91,7 (80,0–97,1) & 84,1 (78,7–86,7) & [0,26; 0,69; 0,66] & 92,5 (80,0–97,1) & 84,6 (79,3–86,0) & [0,21; 0,74; 0,32] & +0,2 \\
R$^3$ — 20-shot† & 90,4 (86,2–95,4) & 68/62 & 92,3 (90,0–95,4) & 84,6 (82,7–86,0) & [0,24; 0,71; 0,66] & 92,5 (90,0–96,2) & 85,1 (81,3–86,7) & [0,23; 0,80; 0,36] & $-$0,3 \\
R$^3$ — 40-shot† & 91,9 (87,3–96,4) & 57/53 & 90,7 (87,3–93,6) & 85,3 (83,3–87,3) & [0,22; 0,71; 0,67] & 90,9 (87,3–92,7) & 85,4 (84,0–86,7) & [0,22; 0,87; 0,28] & $-$0,6 \\
R$^3$ — W ajustado ao rótulo real (teto)† & — & — & — & 86,7 (86,0–87,3) & [0,14; 0,67; 0,73] & — & 86,7 & [0,22; 0,97; 0,00] & — \\
\midrule
R$^4$ — zero-shot† & 50,7 & 149/1 & 72,7 (70,7–74,7) & 63,3 (62,7–64,0) & [0,22; 0,69; 0,10; 0,68] & 77,0 (76,7–77,3) & 67,7 (65,3–70,0) & [0,04; 0,30; 0,12; 0,94] & +6,0 \\
R$^4$ — 4-shot† & 81,4 (68,5–91,8) & 83/63 & 93,5 (89,0–97,9) & 81,5 (76,0–86,7) & [0,22; 0,70; 0,12; 0,63] & 93,4 (87,7–98,6) & 83,0 (79,3–85,3) & [0,25; 0,86; 0,09; 0,14] & +4,3 \\
R$^4$ — 10-shot† & 89,0 (81,4–97,1) & 72/68 & 92,4 (79,3–96,4) & 84,5 (78,0–87,3) & [0,23; 0,72; 0,07; 0,64] & 93,0 (80,0–97,9) & 84,8 (81,3–86,7) & [0,25; 0,95; 0,02; 0,09] & +1,0 \\
R$^4$ — 20-shot† & 90,4 (86,2–95,4) & 68/62 & 93,0 (90,8–96,2) & 85,5 (84,0–88,0) & [0,20; 0,74; 0,07; 0,63] & 92,8 (90,0–96,2) & 85,6 (84,7–88,0) & [0,23; 0,97; 0,03; 0,01] & +0,3 \\
R$^4$ — 40-shot† & 91,9 (87,3–96,4) & 57/53 & 91,7 (89,1–93,6) & 86,1 (84,0–87,3) & [0,18; 0,74; 0,05; 0,65] & 90,9 (88,2–92,7) & 85,4 (84,7–86,7) & [0,21; 0,98; 0,02; 0,00] & +0,4 \\
R$^4$ — W ajustado ao rótulo real (teto)† & — & — & — & 86,2 (85,3–87,3) & [0,14; 0,71; 0,03; 0,69] & — & 86,4 (86,0–86,7) & [0,22; 0,98; 0,00; 0,00] & — \\
\bottomrule
\multicolumn{10}{p{0.95\linewidth}}{\footnotesize Meia-lua (150 pontos, sem split). O LLM vê sempre só x1, x2; R$^3$ acrescenta x1·x2 e R$^4$ acrescenta x1², x2² apenas na estimação da métrica. Colunas e convenções da bloco1\_table\_01: o teto (W ajustado ao rótulo real) só tem acurácia (vs rótulo real), não fidelidade. Em zero-shot (e com nomes de classe neutros) a correspondência entre nome e classe é arbitrária: a acurácia (vs rótulo real) do LLM pode sair abaixo de 50\%; o valor invariante à permutação é max($a$, 1 $-$ $a$). † = em pelo menos uma estimativa o Perceptron devolveu $\hat{W}$ em melhor esforço ($\gamma = 0$, sem margem viável). ``—'' = colapso em classe única (ou quase-colapso: classe minoritária com menos de 5 pontos); na linha colapsada também a acurácia do LLM e a partição ficam como ``—''.} \\
\end{tabular}
\end{table}
